\chapter*{حول مشروع المنطق المفتوح}
\addcontentsline{toc}{chapter}{حول مشروع المنطق المفتوح}


\textit{نص المنطق المفتوح} كتابٌ دراسي تعاوني مفتوح المصدر في المنطق
الفوقي الصوري والمناهج الصورية، يبدأ من مستوى متوسط (أي بعد مقرر
تمهيدي في المنطق الصوري). ومع أنه موجَّه إلى جمهور غير متخصص في
الرياضيات (ولا سيما طلاب الفلسفة وعلوم الحاسوب)، فإنه يتسم بالصرامة.

قد لا تكون تغطية بعض الموضوعات المدرجة حاليًا مكتملة بعد، ولا تزال
أقسام كثيرة بحاجة إلى مراجعة مستفيضة. ونخطط لتوسيع النص مستقبلًا ليشمل
موضوعات أخرى. كما نخطط لإضافة عناصر إلى النص، منها مسرد للمصطلحات،
وقائمة بقراءات إضافية، وملاحظات تاريخية، وصور، وشروح أفضل، وأقسام
تبيّن صلة النتائج بالفلسفة وعلوم الحاسوب والرياضيات، فضلًا عن مزيد من
المسائل والأمثلة. فإذا وجدت خطأً أو كان لديك اقتراح،
\href{https://github.com/OpenLogicProject/OpenLogic/wiki/Contributing}{فيرجى إبلاغ فريق المشروع}.

يعمل المشروع بروح المصدر المفتوح. فليس النص متاحًا بحرية فحسب، بل
نوفر أيضًا مصدره المكتوب بلغة LaTeX بموجب رخصة المشاع الإبداعي لنَسْب
المُصنَّف، وهي تمنح أي شخص الحق في تنزيل عملنا واستعماله وتعديله وإعادة
ترتيبه وتحويله وإعادة توزيعه، شريطة أن ينسب العمل إلى أصحابه على نحو ملائم.
ولمزيد من المعلومات، يرجى زيارة موقع مشروع المنطق المفتوح على
\href{http://openlogicproject.org/}{openlogicproject.org}.
